\BourbakiExercises{Exercises}

\begin{enumerate}
\def\labelenumi{\arabic{enumi})}
\tightlist
\item
  Let \(\mathrm{K}\) be a commutative field of characteristic \(p\) , \(\mathrm{M}\) a \(p\) -primary \(\mathbf{Z}\) -module, \(\mathrm{A}\) the algebra of the group \(\mathrm{M}\) over \(\mathrm{K}\) , and \(\mathfrak{m}\) the set of elements \(\sum_{m} a_{m} e_{m}\) of \(\mathrm{A}\) (cf.~III, §2, No.~6, p.~446) with \(\sum_{m} a_{m} = 0\) .
\end{enumerate}

\begin{enumerate}
\def\labelenumi{\alph{enumi})}
\item
  Prove that A is a local ring, with maximal ideal \(\mathfrak{m}\) , and that \(\mathfrak{m}\) is a nil ideal (cf.~VIII, p.~41, Exercise 10).
\item
  Suppose, moreover, that the homothety \(p_{\mathrm{M}}\) is surjective. Prove that we have \(\mathfrak{m} = \mathfrak{m}^2\) .
\end{enumerate}

\begin{enumerate}
\def\labelenumi{\arabic{enumi})}
\setcounter{enumi}{1}
\item
  Let A be a ring and M an A-module without radical. Prove that the ring of homotheties \(A_{M}\) is without radical (identify \(A_{s}\) with a submodule of \(M^{M}\) ).
\item
  \begin{enumerate}
  \def\labelenumii{\alph{enumii})}
  \tightlist
  \item
    Let A be a ring with semisimple quotient ring \(\mathrm{A} / \Re (\mathrm{A})\) . Prove that we have \(\Re (\mathrm{M}) = \Re (\mathrm{A})\mathrm{M}\) for every A-module M (observe that \(\mathrm{M} / \Re (\mathrm{A})\mathrm{M}\) is an \(\mathrm{A} / \Re (\mathrm{A})\) -module).
  \end{enumerate}
\end{enumerate}

\begin{enumerate}
\def\labelenumi{\alph{enumi})}
\setcounter{enumi}{1}
\tightlist
\item
  Deduce an example of a ring A and a family \((\mathrm{V}_{i})_{i\in\mathrm{I}}\) of A-modules such that \(\Re\big(\prod\mathrm{M}_{i}\big)\neq\prod\Re(\mathrm{M}_{i})\) (use a) for a ring whose radical is not finitely generated).
\end{enumerate}

\begin{enumerate}
\def\labelenumi{\arabic{enumi})}
\setcounter{enumi}{3}
\item
  Let A be an algebra over a commutative field K and a two-sided ideal of A. Let B be the subalgebra of A generated by 1 and a. Prove that the radical of B is contained in that of A (observe that if S is a simple A-module, then we have either ax = 0 for every \(x \in S\) or ax = S for every \(x \neq 0\) in S).
\item
  Let \(A\) be an algebra over a commutative field \(K\) .
\end{enumerate}

\begin{enumerate}
\def\labelenumi{\alph{enumi})}
\tightlist
\item
  Prove that an element of the radical of A that is algebraic over K is nilpotent. b) Suppose that K is an infinite field and that \(\operatorname{Card}(K) > [A : K]\) . Prove that the radical of A is a nil ideal (reason as in Lemma 1 of VIII, p.~47).
\end{enumerate}

\begin{enumerate}
\def\labelenumi{\arabic{enumi})}
\setcounter{enumi}{5}
\tightlist
\item
  Let \(\mathrm{A}\) be a ring.
\end{enumerate}

\begin{enumerate}
\def\labelenumi{\alph{enumi})}
\item
  Let \(\mathfrak{a}\) and \(\mathfrak{b}\) be two-sided ideals of A. If \(\mathfrak{a}\) and \(\mathfrak{b}\) are nilpotent (resp. nil ideals), then so is \(\mathfrak{a} + \mathfrak{b}\) .
\item
  Prove that the sum \(\mathfrak{G}\) of all two-sided nil ideals is a two-sided nil ideal and that the sum \(\mathfrak{N}\) of all nilpotent two-sided ideals is a two-sided nil ideal that contains all nilpotent ideals (VIII, p.~149, Exercise 3).
\item
  Prove that G is the largest semiprime nil ideal (loc. cit.) of A; we say that G is the nilradical of A. If A is commutative, then G and N are both equal to the set of nilpotent elements of A.
\end{enumerate}

¶ 7) Let A be a ring, and let I be the set of ordinals \(\alpha\) such that \(\text{Card}(\alpha) \leqslant \text{Card}(A)\) (Set Theory, III, §6, p.~241, Exercise 10). Define two-sided ideals \(\mathfrak{N}_{\alpha}\) for all \(\alpha \in I\) as follows by transfinite induction. Take \(\mathfrak{N}_0 = \mathfrak{N}\) (Exercise 6, b)). If \(\alpha\) admits a successor \(\alpha + 1\) , then \(\mathfrak{N}_{\alpha+1}\) is the two-sided ideal such that \(\mathfrak{N}_{\alpha+1}/\mathfrak{N}_{\alpha}\) is the sum of the nilpotent two-sided ideals of \(A/\mathfrak{N}_{\alpha}\) . If \(\alpha\) is a limit ordinal, take \(\mathfrak{N}_{\alpha}\) to be the union of the \(N_{\beta}\) for \(\beta < \alpha\) . Let \(\tau\) be the smallest ordinal such that \(N_{\tau+1} = N_{\tau}\) ; the ideal \(P = N_{\tau}\) is called the prime radical of A.

\begin{enumerate}
\def\labelenumi{\alph{enumi})}
\item
  Prove that we have \(N \subset P \subset G\) and that P is the smallest semiprime nil ideal of A. \(^{(1)}\)
\item
  Prove that \(\mathfrak{P}\) is the intersection of the semiprime two-sided ideals of A (cf.~VIII, p.~149, Exercise 3).
\item
  Prove that P is the intersection of the prime two-sided ideals of A (loc. cit., c)). d) Prove that P is the set of elements a of A with the following property: every sequence \((a_{n})_{n\geqslant0}\) of elements of A such that \(a_{0}=a\) and \(a_{n+1}\in a_{n}Aa_{n}\) for every n has finite support (same method as in loc. cit.).
\end{enumerate}

\begin{enumerate}
\def\labelenumi{\arabic{enumi})}
\setcounter{enumi}{7}
\tightlist
\item
  We keep the notation of Exercise 7. Let A be a left Noetherian ring.
\end{enumerate}

\begin{enumerate}
\def\labelenumi{\alph{enumi})}
\item
  Prove that \(\mathfrak{N}\) is the largest nilpotent two-sided ideal of A and that it is equal to the prime radical \(\mathfrak{P}\) of A.
\item
  Prove that every left or right nil ideal of A is nilpotent (apply Exercise 5, a) of VIII, p.~150 to the ring A/N). We therefore have \(\mathfrak{N} = \mathfrak{P} = \mathfrak{G}\) .
\end{enumerate}

\begin{enumerate}
\def\labelenumi{\arabic{enumi})}
\setcounter{enumi}{8}
\item
  Let \(\mathrm{K}\) be a commutative field, \(\mathrm{B}\) the ring \(\mathrm{K}[(\mathrm{X}_n)_{n\geqslant 1}]\) , \(\mathfrak{I}\) the ideal of \(\mathrm{B}\) generated by the sequence \((\mathrm{X}_n^n)\) , and \(\mathrm{A}\) the ring \(\mathrm{B} / \mathfrak{I}\) . Prove that the radical of \(\mathrm{A}\) is a nil ideal but is not nilpotent.
\item
  Let A be a ring.
\end{enumerate}

\begin{enumerate}
\def\labelenumi{\alph{enumi})}
\item
  Let Z be the center of A. Prove that we have \(Z \cap \Re(A) \subset \Re(Z)\) (cf.~Exercise 15, b) below for an example where the inclusion is strict).
\item
  Prove that \(\Re(\mathrm{A})\) contains no nonzero idempotents.
\item
  Prove that the intersection of \(\Re(\mathrm{A})\) and the left socle s of A (VIII, p.~130, Exercise 8) is a two-sided ideal with square zero and that it is the intersection of s and its left annihilator (use Exercises 7 and 8, a) of VIII, p.~130).
\end{enumerate}

\(d\) ) Let \(e\) be an idempotent in A. Prove that the radical of the ring \(e\mathrm{A}e\) is \(e\Re (\mathrm{A})e\) (observe that for \(a\in e\mathrm{A}e\) , we have \((1 - a)(1 - e) = (1 - e)(1 - a) = 1 - e\) ).

\begin{enumerate}
\def\labelenumi{\alph{enumi})}
\setcounter{enumi}{4}
\tightlist
\item
  Prove that the radical of the ring \(\mathbf{M}_{n}(\mathrm{A})\) is the ideal \(\mathbf{M}_{n}(\Re(\mathrm{A}))\) (VIII, p.~114, Example 2).
\end{enumerate}

\begin{enumerate}
\def\labelenumi{\arabic{enumi})}
\setcounter{enumi}{10}
\tightlist
\item
  Let A be a commutative ring and B the ring A{[}X{]}.
\end{enumerate}

\begin{enumerate}
\def\labelenumi{\alph{enumi})}
\item
  First, assume that A has no nonzero nilpotent elements. Let \(f = \sum_{i} a_{i} X^{i}\) and \(g = \sum_{i} b_{i} X^{i}\) be two elements of B. Prove that we have \(a_{p} b_{q} = 0\) for \(p + q > \deg(fg)\) . b) Let N be the nilradical of A (Exercise 6). Prove that a polynomial \(f \in B\) is invertible in B if and only if its constant term is invertible in A and its other coefficients belong to N (apply a) to the image of f in \((\mathrm{A}/\mathfrak{N})[\mathrm{X}]\) .
\item
  Deduce that the radical of B is the ideal \(\Re[X]\) and that it is equal to the nilradical of B.
\end{enumerate}

\begin{enumerate}
\def\labelenumi{\arabic{enumi})}
\setcounter{enumi}{11}
\tightlist
\item
  \begin{enumerate}
  \def\labelenumii{\alph{enumii})}
  \tightlist
  \item
    A ring A is without radical if and only if it is isomorphic to a subring C of a product \(\prod_{i\in I}B_{i}\) of primitive rings such that \(\mathrm{pr}_{i}(\mathrm{C})=\mathrm{B}_{i}\) for every \(i\in I\) (consider the set of classes of simple A-modules).
  \end{enumerate}
\end{enumerate}

\begin{enumerate}
\def\labelenumi{\alph{enumi})}
\setcounter{enumi}{1}
\tightlist
\item
  Let V be an infinite-dimensional vector space over a field D, let B be the ring \(\operatorname{End}_{\mathrm{D}}(\mathrm{V})\) , and let s be its socle (cf.~VIII, p.~130, Exercise 8). Let C be the subring of \(B \times B\) generated by \(s \times s\) and the unit element. Prove that C is without radical but is not isomorphic to a product of primitive rings.
\end{enumerate}

\begin{enumerate}
\def\labelenumi{\arabic{enumi})}
\setcounter{enumi}{12}
\item
  Let \(G\) be a commutative group without torsion and \(K\) a commutative field. Prove that the algebra \(K[G]\) is without radical (endow \(G\) with a total ordering compatible with its group structure, cf.~VI, §1, p.~32, Exercise 20).
\item
  Let \(\mathrm{K}\) be an ordered commutative field and \(\mathrm{G}\) a group. Prove that the algebra of the group \(\mathrm{G}\) over \(\mathrm{K}\) does not admit a nonzero (left or right) nil ideal. (For \(x = \sum a_{g}e_{g}\) , set \(t(x) = a_{1}\) and \(x^{*} = \sum a_{g}e_{g - 1}\) . Observe that if \(x\) is nonzero, then we have \(t((xx^{*})^{2k}) > 0\) for every \(k\geqslant 0\) .) Prove the same property for the algebra of \(\mathrm{G}\) over the field \(\mathrm{K}(i)\) , where \(i^2 = -1\) .
\item
  Let \(V_{0}\) be a vector space of finite dimension n over a field D, and let B be a subring of the ring \(\operatorname{End}_{\mathrm{D}}(\mathrm{V}_{0})\) . Set \(V_{k}=V_{0}\) for \(k\geqslant1\) and \(V=\oplus_{k\geqslant0}V_{k}\) . Consider the set A of endomorphisms u of V for which there exists an integer m (that depends on u) such that we have \(u(V_{k})\subset\oplus_{i\leqslant m}V_{i}\) for \(k\leqslant m\) and \(u(V_{k})\subset V_{k}\) for k\textgreater m and that for k\textgreater m, the endomorphism of \(V_{k}\) obtained by restriction is independent of k and belongs to B.
\end{enumerate}

\begin{enumerate}
\def\labelenumi{\alph{enumi})}
\tightlist
\item
  Prove that A is a primitive ring and that its socle consists of the elements u of A such that \(u(\mathrm{V}_{k}) = 0\) for all but finitely many indices k (VIII, p.~130, Exercise 8).
\item
  Choose suitable D, n, and B to give an example of a primitive ring that has a quotient ring with nilpotent nonzero radical and an example of a primitive ring whose center has nonzero radical.
\end{enumerate}

\begin{enumerate}
\def\labelenumi{\arabic{enumi})}
\setcounter{enumi}{15}
\tightlist
\item
  Let A be a ring and \(\mathfrak{a}\) a two-sided ideal contained in the radical of A. Let M and N be A-modules and \(u: \mathrm{M} \to \mathrm{N}\) a homomorphism. We denote by \(\overline{u}: \mathrm{M} / \mathfrak{aM} \to \mathrm{N} / \mathfrak{aN}\) the homomorphism of A/ \(\mathfrak{a}\) -modules deduced from \(u\) .
\end{enumerate}

\begin{enumerate}
\def\labelenumi{\alph{enumi})}
\item
  Suppose that M and N are finitely generated and N is projective. If \(\overline{u}\) is bijective, then so is u.
\item
  Give an example with M and N free of rank 1 and \(\overline{u}\) injective but \(\operatorname{Ker}(u) \neq 0\) . c) Suppose that M and N are projective. If \(\overline{u}\) induces an isomorphism to a direct factor, then u is injective (reduce to the case when M and N are free and equal and \(\overline{u}\) is the identity, and then use a)). If, moreover, M is finitely generated, then the image of u is a direct factor of N.
\item
  Suppose that M is finitely generated and N is projective. If \(\overline{u}\) is bijective, then so is u (apply c) to a homomorphism \(v: N \to M\) such that \(\overline{v} = \overline{u}^{-1}\) to conclude that N is finitely generated).
\item
  If M is projective and M/ \(\mathfrak{aM}\) is zero, then M is zero.
\end{enumerate}

\begin{enumerate}
\def\labelenumi{\arabic{enumi})}
\setcounter{enumi}{16}
\tightlist
\item
  Let \(\mathbf{K}\) be a commutative field and \(\mathbf{L}\) a set.
\end{enumerate}

\begin{enumerate}
\def\labelenumi{\alph{enumi})}
\tightlist
\item
  Prove that there exists a K-algebra A admitting a basis consisting of the unit element, a family \((c_{\lambda})_{\lambda \in \mathrm{L}}\) , and a family \((r_{\lambda \mu})_{(\lambda ,\mu)\in \mathrm{L}\times \mathrm{L}}\) , with multiplication table
\end{enumerate}

\[
\begin{array}{r} c _ {\lambda} ^ {2} = c _ {\lambda}, \qquad c _ {\lambda} c _ {\mu} = r _ {\lambda \mu} \quad \mathrm{if} \lambda \neq \mu \\ c _ {\lambda} r _ {\mu \nu} = \delta_ {\lambda \mu} r _ {\mu \nu}, \qquad r _ {\lambda \mu} c _ {\nu} = \delta_ {\mu \nu} r _ {\lambda \mu} \\ r _ {\lambda \mu} r _ {\nu \rho} = 0, \end{array}
\]

where \(\delta_{\mu\nu}\) is the Kronecker delta function. Prove that \(\Re(\mathrm{A})\) is the subalgebra of A generated by the \(r_{\lambda\mu}\) ; we have \(\Re(\mathrm{A})^{2}=0\) .

\begin{enumerate}
\def\labelenumi{\alph{enumi})}
\setcounter{enumi}{1}
\tightlist
\item
  Suppose \(\operatorname{Card}(\mathrm{L}) \geqslant 2\) . Prove that the center Z of A is the K-vector space generated by 1, the \(r_{\lambda\lambda}\) for \(\lambda \in L\) , and, if L is finite, the element \(\sum_{\lambda} c_{\lambda} - \sum_{\mu \neq \nu} r_{\mu\nu}\) . Deduce that for \(\lambda \in L\) , there are no idempotents in Z congruent to \(c_{\lambda} \pmod{\Re(A)}\) .
\end{enumerate}

¶ 18) a) Let A be a ring, a a two-sided nil ideal of A, and \((\overline{e}_{n})_{n\geqslant1}\) an orthogonal sequence of idempotents in A/a. Prove that there exists an orthogonal sequence \((e_{n})_{n\geqslant1}\) of idempotents in A such that \(\overline{e}_{n}\) is the class of \(e_{n}\) in A/a for every n.~(Let \(\overline{e}\) and \(\overline{e}'\) be two orthogonal idempotents in A/a, \(e'\) an idempotent in A lifting \(\overline{e}'\) , and let a be a representative of e in A. Observe that we can modify a to have \(ae' = e'a\) , and then, as in Proposition 7 of VIII, p.~159, to have \(a^{2} = a\) .)

\begin{enumerate}
\def\labelenumi{\alph{enumi})}
\setcounter{enumi}{1}
\tightlist
\item
  Let K be a commutative field, L an uncountable set, and A the K-algebra defined in the previous exercise. Prove that there does not exist any orthogonal family \((e_{\lambda})_{\lambda\in\mathrm{L}}\) of idempotents in A satisfying \(e_{\lambda}\equiv c_{\lambda} \pmod{\Re(\mathrm{A})}\) for every \(\lambda\in L\) . (Consider the set H of elements \((\lambda,\mu)\) of \(L\times L\) such that \(c_{\lambda}(e_{\mu}-c_{\mu})=0\) . Observe that the sets \(H\cap(\{\lambda\}\times L)\) are finite for every \(\lambda\in L\) and the sets \(\mathbb{C}H\cap(L\times\{\mu\})\) are finite for every \(\mu\in L\) .)
\end{enumerate}

\begin{enumerate}
\def\labelenumi{\arabic{enumi})}
\setcounter{enumi}{18}
\item
  Let A be a ring, a a two-sided ideal of A contained in \(\Re(\mathrm{A})\) , let e and \(e'\) be idempotents in A, and let \(\overline{e}\) and \(\overline{e}'\) their classes in the ring \(\overline{A} = A/a\) . Prove that if the idempotents \(\overline{e}\) and \(\overline{e}'\) are equivalent (VIII, p.~39, Exercise 5), then so are e and \(e'\) (observe that Ae is a projective cover of the A-module \(\overline{A}\overline{e}\) ).
\item
  Let \(\mathrm{A}\) be an algebra over a commutative field \(\mathrm{K}\) such that \(\Re(\mathrm{A})\) is a nil ideal and that the algebra \(\mathrm{A} / \Re(\mathrm{A})\) is isomorphic to a finite product of matrix algebras over \(\mathrm{K}\) . Prove that there exists a subalgebra \(\mathrm{S}\) of \(\mathrm{A}\) such that the \(\mathrm{K}\) -vector space \(\mathrm{A}\) is the direct sum of \(\mathrm{S}\) and \(\Re(\mathrm{A})\) (use Exercises 18 and 19 above, as well as Exercise 19 of VIII, p.~94).
\item
  Let A be a ring, \(u: P \to M\) a surjective homomorphism of A-modules, and N its kernel. We denote by \(\operatorname{End}_{\mathrm{A}}(\mathrm{P}; \mathrm{N})\) (resp. \(\operatorname{Aut}_{\mathrm{A}}(\mathrm{P}; \mathrm{N})\) ) the ring of endomorphisms (resp. the group of automorphisms) u of P such that \(u(\mathrm{N}) \subset \mathrm{N}\) .
\end{enumerate}

\begin{enumerate}
\def\labelenumi{\alph{enumi})}
\tightlist
\item
  Define an exact sequence
\end{enumerate}

\[
0 \to \mathrm{Hom} _ {\mathrm{A}} (\mathrm{P}, \mathrm{N}) \xrightarrow {i} \mathrm{End} _ {\mathrm{A}} (\mathrm{P}; \mathrm{N}) \xrightarrow {p} \mathrm{End} _ {\mathrm{A}} (\mathrm{M})
\]

where the image of \(i\) is a two-sided ideal \(\mathfrak{I}\) of the ring \(\operatorname{End}_{\mathrm{A}}(\mathrm{P};\mathrm{N})\) .

\begin{enumerate}
\def\labelenumi{\alph{enumi})}
\setcounter{enumi}{1}
\tightlist
\item
  Suppose that \((P, u)\) is a projective cover of M. Then p is surjective and induces a surjective homomorphism \(p' : \operatorname{Aut}_{\mathrm{A}}(\mathrm{P}; \mathrm{N}) \to \operatorname{Aut}_{\mathrm{A}}(\mathrm{M})\) . We have an exact sequence of groups
\end{enumerate}

\[
0 \to 1 _ {\mathrm{P}} + \Im \to \mathrm{Aut} _ {\mathrm{A}} (\mathrm{P}; \mathrm{N}) \xrightarrow {p ^ {\prime}} \mathrm{Aut} _ {\mathrm{A}} (\mathrm{M}) \to 0;
\]

if P is free and N is not reduced to 0, then the homomorphism \(p'\) is not injective.

\begin{enumerate}
\def\labelenumi{\arabic{enumi})}
\setcounter{enumi}{21}
\tightlist
\item
  Let \(A\) be a commutative integral domain that has only a finite number \(s > 1\) of maximal ideals \(\mathfrak{m}_1, \ldots, \mathfrak{m}_s\) .
\end{enumerate}

\begin{enumerate}
\def\labelenumi{\alph{enumi})}
\item
  Prove that the ring A/ \(\Re(A)\) is isomorphic to the product of the fields A/mi.
\item
  Prove that the A-modules \(A/m_{i}\) do not admit a projective cover even though they are projective and finitely generated as \(A/\Re(A)\) -modules (observe that A is a projective cover of \(A/\Re(A)\) ).
\item
  Let K be a commutative field. Prove that the subring A of the field \(\mathrm{K(T)}\) consisting of the rational fractions that can be written as \(\mathrm{P(T)/Q(T)}\) with \(\mathrm{Q(0) \neq 0}\) and \(\mathrm{Q(1) \neq 0}\) has the required properties.
\end{enumerate}

\begin{enumerate}
\def\labelenumi{\arabic{enumi})}
\setcounter{enumi}{22}
\item
  Let A be a ring and a a (left) ideal of A. Suppose that for every sequence \((a_{n})_{n\geqslant1}\) of elements of a, there exists an integer m such that \(a_{1}\cdots a_{m}=0\) . Prove that we have \(aM\neq M\) for every nonzero A-module M (if not, construct, by induction, sequences \((a_{n})_{n\geqslant1}\) in a and \((x_{n})_{n\geqslant1}\) in M such that we have \(a_{1}\cdots a_{p}x_{p}\neq0\) for every p).
\item
  Let \(A\) be a ring and \(\mathfrak{r}\) its radical. Prove that the following properties are equivalent:
\end{enumerate}

\begin{enumerate}
\def\labelenumi{(\roman{enumi})}
\tightlist
\item
  Every finitely generated left A-module admits a projective cover.
\end{enumerate}

(i') Every finitely generated right A-module admits a projective cover.

\begin{enumerate}
\def\labelenumi{(\roman{enumi})}
\setcounter{enumi}{1}
\tightlist
\item
  The ring \(\mathrm{A} / \mathfrak{r}\) is semisimple, and every idempotent in \(\mathrm{A} / \mathfrak{r}\) is the class of an idempotent in \(\mathrm{A}\) .
\end{enumerate}

(To prove (i)⇒(ii), observe that every decomposition \(A/r = a_{1} \oplus a_{2}\) into a direct sum of ideals lifts to a decomposition \(A = P_{1} \oplus P_{2}\) , where \(P_{i}\) is a projective cover of \(a_{i}\) . Show that, moreover, (ii) implies that every A/r-module is of the form \((\mathrm{A}/\mathfrak{r}) \otimes_{\mathrm{A}} \mathrm{P}\) , where P is a projective A-module.)

¶ 25) Let A be a ring and r its radical. Prove that the following conditions are equivalent:

\begin{enumerate}
\def\labelenumi{(\roman{enumi})}
\tightlist
\item
  Every left A-module admits a projective cover.
\end{enumerate}

(i') Every right A-module admits a projective cover.

\begin{enumerate}
\def\labelenumi{(\roman{enumi})}
\setcounter{enumi}{1}
\tightlist
\item
  The ring A/r is semisimple, and for every sequence \((a_{n})_{n\geqslant1}\) of elements of r, there exists an integer m such that \(a_{1}\cdots a_{m}=0\) .
\end{enumerate}

(Prove (ii)⇒(i) in the same way as the corresponding implication of Exercise 24, using Exercise 23. Assume that condition (i) is satisfied. Consider the submodule M of A \(^{(N)}\) generated by \(e_1 - a_1e_2\) , \(e_2 - a_2e_3\) , \ldots{} Deduce from Exercise 16, e) applied to a projective cover of A \(^{(N)}\) /M that we have M = A \(^{(N)}\) , which implies that \(a_1 \cdots a_n = 0\) for n sufficiently large.)

\begin{enumerate}
\def\labelenumi{\arabic{enumi})}
\setcounter{enumi}{25}
\tightlist
\item
  A ring A is called a Zorn ring if every left ideal of A that is not a nil ideal contains a nonzero idempotent.
\end{enumerate}

\begin{enumerate}
\def\labelenumi{\alph{enumi})}
\item
  Prove that in a Zorn ring, every right ideal that is not a nil ideal contains a nonzero idempotent (observe that for every nonnilpotent element \(a\) of A, there exists an \(x \in A\) such that \(xaxa = xa \neq 0\) ).
\item
  Prove that the radical of a Zorn ring is a nil ideal and that a Zorn ring without zero divisors is a field.
\item
  Prove that a primitive ring with nonzero socle is a Zorn ring (cf.~VIII, p.~129, Exercise 4, b) and p.~130, Exercise 7, a)). Give an example of a primitive ring that is not a Zorn ring (cf.~VIII, p.~132, Exercise 13) and an example of a primitive ring whose socle is nonzero but whose center is not a Zorn ring (use the method of Exercise 12).
\end{enumerate}

\(d\) ) Let K be a commutative field and A the subring of \(\mathrm{K}(\mathrm{X})^{\mathbf{N}}\) consisting of the sequences \(s = (f_n)\) for which there exists a polynomial \(\varphi(s) \in \mathrm{K}[\mathrm{X}]\) such that \(f_n\) is equal to \(\varphi(s)\) for \(n\) sufficiently large. Prove that A is a commutative Zorn ring, without radical, containing elements that are neither invertible nor zero divisors, but that its image \(\varphi(A)\) is not a Zorn ring.

¶ 27) A ring A is called absolutely flat, or von Neumann regular, if for every element a of A, there exists an \(x \in A\) such that axa = a.

\begin{enumerate}
\def\labelenumi{\alph{enumi})}
\item
  A ring A is absolutely flat if and only if every monogenous left (or right) ideal of A admits a supplement in \(A_{s}\) (that is, is generated by an idempotent).
\item
  Show that every finitely generated ideal of an absolutely flat ring is generated by an idempotent. (Reduce to the case of an ideal \(Ae + Af\) , where e and f are idempotents. Then, by replacing Ae with \(Ae(1 - f)\) , reduce to the case when ef = 0, and consider the element \(e + f - fe\) .) A (left or right) Noetherian absolutely flat ring is semisimple.
\item
  Deduce from b) that the intersection of two monogenous left (resp. right) ideals of an absolutely flat ring is monogenous (consider the right annihilators).
\item
  Every quotient ring of an absolutely flat ring is absolutely flat, as is every product of absolutely flat rings.
\item
  Prove that the radical of an absolutely flat ring is reduced to 0. Deduce that every two-sided ideal of A is the intersection of the maximal left ideals containing it.
\end{enumerate}

\(f\) ) The center of an absolutely flat ring is absolutely flat (prove that if \(a \in \mathbb{Z}\) and \(x \in \mathrm{A}\) satisfy \(a^2 x = a\) , then we have \(a^2 x^n \in \mathbb{Z}\) for every \(n \geqslant 0\) and \(a(a^2 x^3)a = a\) ). \(g\) ) Prove that the endomorphism ring of a vector space is an absolutely flat ring.

¶ 28) Let A be an absolutely flat ring (Exercise 27) and n an integer.

\begin{enumerate}
\def\labelenumi{\alph{enumi})}
\item
  Prove that every finitely generated submodule of \(A^{n}\) admits a supplementary submodule. (Reason by induction on n, by showing that \(M \cap A^{n-1}\) is finitely generated and choosing supplementary submodules of \(M \cap A^{n-1}\) in \(A^{n-1}\) and of the image of M in \(Ae_{n}\) .)
\item
  Deduce that the ring \(\mathbf{M}_n(\mathrm{A})\) is absolutely flat (cf.~VIII, p.~114, Example 2).
\end{enumerate}

\begin{enumerate}
\def\labelenumi{\arabic{enumi})}
\setcounter{enumi}{28}
\item
  Let V be a vector space of countable dimension over a commutative field K, and let \(A = \text{End}_{\text{K}}(\text{V})\) . Prove that there exists a unique maximal two-sided ideal that is \(\text{End}_{\text{K}}^{\text{f}}(\text{V})\) , the set of endomorphisms of finite rank (cf.~Exercise 2 of VIII, p.~128), but that the radical of A is 0.
\item
  \begin{enumerate}
  \def\labelenumii{\alph{enumii})}
  \tightlist
  \item
    Let A be a ring that contains no nilpotent elements other than 0. Prove that every idempotent e in A belongs to the center of A (consider the elements \((1-e)xe\) and \(ex(1-e)\) ).
  \end{enumerate}
\end{enumerate}

\begin{enumerate}
\def\labelenumi{\alph{enumi})}
\setcounter{enumi}{1}
\item
  Let A be an absolutely flat ring that contains no nilpotent elements other than 0, and let a be a nonzero element of A. Prove that there exists an element x of A such that e = ax = xa is idempotent and ea = ae = a. Deduce that for every \(y \in A\) , there exists a \(z \in A\) such that ya = az and, consequently, that every left or right ideal of A is two-sided.
\item
  Prove that no quotient ring of A contains a nonzero nilpotent element (use b)). Deduce that A is isomorphic to a subring C of a product \(\prod_{\iota\in I}D_{\iota}\) of fields such that \(\mathrm{pr}_{\iota}(\mathrm{C})=\mathrm{D}_{\iota}\) for every \(\iota\in I\) (deduce from b) and Exercise 4 of VIII, p.~129 that an absolutely flat primitive ring without nilpotent elements is a field, and use Exercise 12 above).
\end{enumerate}

\(d\) ) Let \((\mathrm{D}_{\iota})_{\iota \in \mathbb{I}}\) be an infinite family of fields and B the ring \(\prod_{\iota \in \mathrm{I}}\mathrm{D}_{\iota}\) . For any subset H of I, denote by \(\mathrm{B_H}\) the two-sided ideal of B consisting of the families \((b_{\iota})_{\iota \in \mathrm{I}}\) such that \(b_{\iota} = 0\) for \(\iota \notin \mathrm{H}\) . Let \(\mathfrak{F}\) be an increasing directed set of subsets of I forming a cover of I, and let A be the subring of B generated by 1 and the \(\mathrm{B_H}\) for \(\mathrm{H}\in \mathfrak{F}\) . Prove that A is absolutely flat and contains no nonzero nilpotent elements.

\begin{enumerate}
\def\labelenumi{\arabic{enumi})}
\setcounter{enumi}{30}
\tightlist
\item
  Let K be a commutative field, and let A be the set of matrices of the form \(\begin{pmatrix} a & b \\ 0 & c \end{pmatrix}\) with \(a, b, c \in K\) . Verify that A is a subring of the matrix ring and that its radical is determined by a = c = 0. Prove that the image of \(\begin{pmatrix} 1 & 0 \\ 0 & 0 \end{pmatrix}\) in A/Re(A) is an idempotent element of the center of A/Re(A).
\end{enumerate}

